\subsection{The source's weight-zero constancy argument}
\label{sec:ftb-key-weight-zero}

This subsection supplies the weight-zero statement used to normalize
the holomorphic key invariant.  It follows the frozen source's
Fourier estimate and rank-one modular curves.  The sole
one-dimensional modular input is the usual theorem that a
holomorphic modular form of weight zero for
\(\mathrm{SL}_2(\mathbb Z)\), bounded at its cusp, is constant.
It follows, for example, by the removable-singularity theorem and
the maximum principle on the compactified standard fundamental
domain.  No higher-dimensional compactification or boundary
extension theorem is used here.

\begin{srlemma}[Fourier bound on a Siegel cone]
\label{lem:ftb-ph-koecher}
Let \(g\geq2\), \(N\geq1\), and let \(F\) be holomorphic on
\(\mathfrak H_g\).  Suppose
\[
 \begin{aligned}
 F(\tau+NB)&=F(\tau) &&(B=B^t\in M_g(\mathbb Z)),\\
 F(U\tau U^t)&=F(\tau)
 &&(U\in\mathrm{SL}_g(\mathbb Z),\ U\equiv1\bmod N).
 \end{aligned}                                             \tag{W1}\label{eq:ftb-ph-periods}
\]
For every \(c>0\) there is a finite real \(K_c\) such that
\[
       |F(\tau)|\leq K_c\quad
          \text{whenever }\operatorname{Im}\tau\geq c1.
 \tag{W2}\label{eq:ftb-ph-cone-bound}
\]
\end{srlemma}
\begin{proof}
For clarity we give the source's full-matrix Fourier construction;
this avoids off-diagonal half-integer conventions.
Put
\[
 \Phi(z)=N(z+z^t),\qquad
 \mathcal U=\{z\in M_g(\mathbb C):\operatorname{Im}\Phi(z)>0\},
 \qquad f(z)=F(\Phi(z)).
\]
This is an open convex tube.  The function \(f\) has period \(1\)
in each of its \(g^2\) coordinates, by \eqref{eq:ftb-ph-periods}.
For \(m\in M_g(\mathbb Z)\), the integral
\[
 c_m=\int_{[0,1]^{g^2}}
 f(x+z)e^{-2\pi i\langle m,x+z\rangle}\,dx
 \tag{W3}\label{eq:ftb-tube-coeff}
\]
is independent of \(z\in\mathcal U\).
Indeed the integrand is holomorphic and has unit periods.
Its cube integral is holomorphic in \(z\), by differentiation
under an integral over a compact cube, and invariant under every
real translation, by the periodic cube-change formula.
Its real derivatives vanish; complex differentiability therefore
makes all its complex derivatives vanish.  Convexity makes it
constant.  Here \(\langle m,z\rangle=\sum_{ij}m_{ij}z_{ij}\).

Set \(a_m=|c_m|\).
If a real imaginary slice is bounded by \(M\), its Fourier
coefficient estimate gives
\[
          a_m e^{-2\pi\langle m,Y\rangle}\leq M.
 \tag{W4}\label{eq:ftb-tube-coeff-bound}
\]
At a fixed \(z\), any summable nonnegative majorant \(b_m\) of
the left side gives \(|f(z)|\leq\sum_m b_m\).
To justify this last assertion without assuming pointwise Fourier
convergence, regard \(x\mapsto f(x+z)\) as a continuous function
on the unit real torus.  Its Fourier coefficients are the
corresponding expressions in \eqref{eq:ftb-tube-coeff}.
An absolutely summable coefficient sequence has a uniformly
convergent Fourier series.  Fourier uniqueness, or the elementary
approximate-identity argument on that torus, identifies its sum
with the original continuous function.  Taking norms proves the
assertion.

The unit periods make \(F\) bounded on each fixed positive
imaginary slice: reduce the real symmetric entries to a compact
cube.  Fix \(c>0\), and let \(M\geq0\) bound the slice
\((c/2)1\).  Set \(s=c/(4N)>0\).
For every antisymmetric real \(A\), the slice \(Y=s1+A\)
of \(\mathcal U\) has
\(N(Y+Y^t)=(c/2)1\); hence
\[
 a_m e^{-2\pi\langle m,s1+A\rangle}\leq M.
 \tag{W5}\label{eq:ftb-anti-test}
\]
For every \(U\in\mathrm{SL}_g(\mathbb Z)\) congruent to \(1\bmod N\),
invariance and conjugation by the integral inverse of \(U\)
give the same bound on the slice \((c/2)UU^t\).  Thus
\[
 a_m e^{-2\pi s\,\operatorname{tr}(U^tmU)}\leq M.
 \tag{W6}\label{eq:ftb-congruence-test}
\]

If \(a_m>0\), \eqref{eq:ftb-anti-test} forces \(m=m^t\):
otherwise choose \(A=t(E_{ij}-E_{ji})\) and let
\(\langle m,A\rangle\) tend to \(-\infty\).
Next \(m\) is positive semidefinite.
If not, continuity gives a rational vector \(x\) with \(x^tmx<0\);
clear denominators to take \(x\in\mathbb Z^g\).
Choose \(i\) with \(x_i\ne0\), choose \(j\ne i\), and put
\(y=x_je_i-x_ie_j\ne0\).
Then \(y^tx=0\) and
\[
 U_k=1+Nkxy^t\in\mathrm{SL}_g(\mathbb Z),\quad
 U_k\equiv1\bmod N,\quad
 \operatorname{tr}(U_k^tmU_k)
 =\operatorname{tr}m+2Nk\,x^tmy
                 +(Nk)^2(x^tmx)\,y^ty .
\]
The determinant formula for a rank-one perturbation proves the
first assertion.  The last polynomial tends to \(-\infty\).
It contradicts \eqref{eq:ftb-congruence-test}, because \(a_m>0\).
Rational vectors are dense, so this proves real positive
semidefiniteness as well.

For such an \(m\), put \(t=\operatorname{tr}m\geq0\).
The inequalities obtained by evaluating its quadratic form at
\(e_i\pm e_j\) give \(|m_{ij}|\leq t\).  Consequently
\[
       \sum_{ij}|m_{ij}|\leq g^2t,\qquad
       \langle m,\operatorname{Im}\tau\rangle\geq ct
                    \quad(\operatorname{Im}\tau\geq c1).
\]
Use \eqref{eq:ftb-tube-coeff-bound} at \(Y=s1\).
With \(\varepsilon=\pi c/(2Ng^2)>0\) the preceding estimates give
\[
 a_m e^{-(\pi/N)\langle m,\operatorname{Im}\tau\rangle}
       \leq M\exp\!\left(-\varepsilon\sum_{ij}|m_{ij}|\right).
 \tag{W7}\label{eq:ftb-ph-summable}
\]
The same inequality is trivial if \(a_m=0\).
The right side is summable over \(M_g(\mathbb Z)\): it is a
finite product of two-sided geometric series.
Apply the torus Fourier bound at \(z=\tau/(2N)\).
This proves \eqref{eq:ftb-ph-cone-bound}, with exactly the
source's summable majorant.
\end{proof}

\begin{srlemma}[A finite modular family is constant]
\label{lem:ftb-finite-elliptic-family}
Let \(\mathcal F\) be a finite set of holomorphic functions on
the upper half-plane, each bounded as the imaginary part tends
to infinity.  Suppose precomposition by \(S:z\mapsto-1/z\)
and \(T:z\mapsto z+1\) takes \(\mathcal F\) into itself.
Then every function in \(\mathcal F\) is constant.
\end{srlemma}
\begin{proof}
Precomposition is injective, since \(S,T\) are bijections, so it
permutes the finite set.
For each \(w\in\mathbb C\),
\(\prod_{f\in\mathcal F}(w-f(z))\) is therefore a holomorphic
weight-zero modular function for the full modular group, bounded
at infinity and hence at all its equivalent cusps.
The one-dimensional modular input makes it constant.
Fix \(z_0\).  Taking \(w=f(z)\) in that polynomial identity shows
that \(f(z)\) belongs to the finite set
\(\{h(z_0):h\in\mathcal F\}\).
The upper half-plane is connected; continuity makes \(f\)
constant.
\end{proof}

\begin{srproposition}[Holomorphic symplectic invariants have weight-zero constancy]
\label{prop:ftb-ph-holo}
If \(g\geq2\) and \(H:\mathfrak H_g\to\mathbb C\) is holomorphic
and invariant under \(\mathrm{Sp}_{2g}(\mathbb Z)\), then \(H\)
is constant.
\end{srproposition}
\begin{proof}
The symplectic translations and diagonal symplectic matrices give
\eqref{eq:ftb-ph-periods} with \(N=1\).
Lemma~\ref{lem:ftb-ph-koecher} gives a bound on every cone
\(\operatorname{Im}\tau\geq c1\).
We use this bound on the following explicit rank-one curves.

Fix a coordinate \(e\), write \(\epsilon=e_e\), and suppose
\(Z=Z^t\) has zero \(e\)-th column and
\[
       \operatorname{Im}Z+yuu^t>0
                    \quad\text{for all }u_e=1,\ y>0.
\]
For real symmetric \(R\) with \(R_{ee}=0\), and real \(u_e=1\),
put \(C_{R,u}(z)=Z+R+zuu^t\).
For a fixed integer \(N\geq1\), consider all functions
\(z\mapsto H(C_{R,u}(z))\) satisfying
\[
    Nu_i\in\mathbb Z,\quad NR_{ie}\in\mathbb Z,\quad
                             N^2R_{ij}\in\mathbb Z.
 \tag{W8}\label{eq:ftb-curve-denominators}
\]
There are only finitely many such \emph{functions}.
To see this, set \(k_e=0,\ k_i=\lfloor u_i\rfloor\) for \(i\ne e\).
Conjugation by the integral unimodular rank-one shear
\(1-k\epsilon^t\) changes the curve parameters to
\[
 u_1=u-k,\qquad
 R_1=R-k(Re)^t-(Re)k^t .
\]
It fixes \(Z\), since the \(e\)-th row and column of \(Z\) are zero.
Subtract the symmetric integral matrix of floors of \(R_1\);
symplectic translation changes no value of \(H\).
Now \(u_1\) and the fractional entries of \(R_1\) lie in
\([-1,1]\), and their denominators divide \(N^2\).
There are finitely many choices.

Each function is holomorphic and bounded at infinity.
For the bound, its imaginary part for \(y\geq1\) is at least
\(\operatorname{Im}Z+uu^t>0\); choose \(c1\) below this fixed
positive matrix and apply the cone bound.
The finite family is stable under \(T\), which changes the
parameters, up to the integral \(\epsilon\epsilon^t\) translation,
to
\[
                         (R+uu^t-\epsilon\epsilon^t,\ u).
\]
It is also stable under \(S\).  The symplectic partial inversion
in coordinate \(e\) gives the parameters
\[
 u'=\epsilon+Re,\qquad
 R'=R-u(\epsilon+Re)^t-(\epsilon+Re)u^t
                                      +2\epsilon\epsilon^t.
 \tag{W9}\label{eq:ftb-curve-S}
\]
These have \(u'_e=1,\ R'_{ee}=0\), and satisfy
\eqref{eq:ftb-curve-denominators}: multiply each formula by \(N\)
or \(N^2\) and use the displayed integrality conditions.
The partial-inversion identity is checked entry by entry after
using \(Z e=0,\ u_e=1,\ R_{ee}=0\); it asserts that
\(H(C_{R,u}(-1/z))=H(C_{R',u'}(z))\).
Lemma~\ref{lem:ftb-finite-elliptic-family} makes every such
rational-parameter curve constant.

For real \(u_e=1\) and \(v_e=0\), take
\(R=\epsilon v^t+v\epsilon^t\).
Approximate every coordinate of \(u,v\) by
\(\lfloor(n+1)x\rfloor/(n+1)\).
The \(e\)-coordinates remain exactly \(1,0\), respectively, and
the denominator conditions hold with \(N=n+1\).
Passing to the limit by continuity gives constancy on this
real-parameter curve as well.

We now exhibit such a curve through every \(\tau=X+iY\).
Set
\[
 h=Y_{ee},\quad u=Ye/h,\quad
 v=Xe-X_{ee}u,\quad
 R=\epsilon v^t+v\epsilon^t,\quad
 Z=\tau-R-\tau_{ee}uu^t .
\]
Then \(h>0,\ u_e=1,\ v_e=0,\ Ze=0\) and
\(\operatorname{Im}Z=Y-(Ye)(Ye)^t/h\).
For every \(u'_e=1,\ y>0\),
\[
                   Y-(Ye)(Ye)^t/h+yu'(u')^t>0.
\]
Indeed for \(x\ne0\) put
\(x'=x-(x^tYe/h)e\).
Its quadratic value is
\((x')^tYx'+y(x^tu')^2\).
If \(x'\ne0\), the first term is positive.
If \(x'=0\), \(x\) is a nonzero multiple of \(e\) and the second
term is positive.  Thus the curve has the required domain and
passes through \(\tau\) at \(z=\tau_{ee}\).
Its constancy gives
\[
                     dH_\tau((Ye)(Ye)^t)=0.
\]
Conjugate this assertion by integral unimodular matrices.
Their \(e\)-th rows range over the primitive integer vectors,
so it gives
\(dH_\tau((Yv)(Yv)^t)=0\) for every such \(v\).
The vectors \(e_i\) and \(e_i+e_j\), \(i\ne j\), are primitive
rows of unimodular matrices; their outer squares span all complex
symmetric matrices, by subtracting the diagonal squares.
Since \(Y\) is invertible, \(dH_\tau\) vanishes on the whole
symmetric tangent space.
Finally \(\mathfrak H_g\) is convex.  Integrating along its line
segments proves that \(H\) is constant.
\end{proof}

\begin{srcorollary}[Finite symplectic families]
\label{cor:ftb-finite-sp-family}
Let \(g\geq2\).  A finite collection of holomorphic functions on
\(\mathfrak H_g\), permuted by precomposition by every integral
symplectic transformation, consists of constant functions.
\end{srcorollary}
\begin{proof}
Discard duplicate functions.
For each \(w\in\mathbb C\) the polynomial
\(\prod_f(w-f(\tau))\) is holomorphic and symplectic invariant,
because precomposition is an injective map of this finite set
and hence a permutation.
Proposition~\ref{prop:ftb-ph-holo} makes it constant.
Its roots at a fixed period form a finite set containing every
value of every \(f\).  Connectedness and continuity conclude
exactly as in Lemma~\ref{lem:ftb-finite-elliptic-family}.
\end{proof}

\paragraph{Application to the powered key rule.}
In this paragraph assume explicitly \(g\geq2\).
For a field-functorial section-independent rule satisfying
\eqref{eq:ftb-koszul-homog}, the period construction gives one
holomorphic value function \(U_{a,b}\) for each half-characteristic.
The theta transformation formula permutes this finite family.
More precisely, for \(\delta\in\mathrm{Sp}_{2g}(\mathbb Z)\),
write \(j=C\tau+D\).  The theta formula gives a nonzero multiplier
\(r\) with \(r^8=(\det j)^4\); reparametrizing the uniformization
changes the normal-form data to
\[
        \tau'=\delta\tau,\qquad M'=Mj^{-1},\qquad c'=c/r.
\]
It preserves
\[
                         \frac{(\det M')^4}{(c')^8}
                             =\frac{(\det M)^4}{c^8}.
 \tag{W10}\label{eq:ftb-normal-ratio-sp}
\]
The Gaussian exponential in the theta formula is incorporated
in the reparametrized line frame, not discarded.
The homogeneity law and \eqref{eq:ftb-normal-ratio-sp} show that
the normal value functions are permuted, with no scalar factor.
Apply Corollary~\ref{cor:ftb-finite-sp-family}.
For the zero characteristic there is a constant \(a_0\) such that
\[
             \kappa(x)=a_0
                    \left(\frac{(\det M)^4}{c_0^8}\right)^c
 \tag{W11}\label{eq:ftb-key-unscaled-constant}
\]
whenever \(x\) is in zero-characteristic normal form
\((\tau,M,c_0)\).
To pass from the normalized \(M=1,c_0=1\) case to this formula,
replace the theta section by \(c_0^{-1}\vartheta\) and the Lie
basis by \(M^{-1}v\), then apply
\eqref{eq:ftb-koszul-homog}.
The symplectic transformation formula here is the classical
Poisson theta transformation with its eighth-power identity;
its source implementation uses the squared formula and the
identity theorem to choose \(r\), and then changes the actual
uniformization frame.

This proves constancy, not yet the rationality or integral unit
property of \(a_0\), and not yet an exponent-one key rule.
Those are supplied by the subsequent arithmetic normalization
and connected-family descent.

\paragraph{Frozen source concordance.}
All proof paths below have prefix
\srcid{src/Development/AbelianVarietyTheory/Proofs/ThetaLawProofs/}.
The tube Fourier construction and cone estimate are
\srcid{NormLeTsumOfCubeFourierCoeffLe.lean}
(\srcid{958fca26bfcc625698732b1a0dafee0d51889f7b});
its explicit congruence transvection is
\srcid{ExistsNatMulQuadraticLt.lean}
(\srcid{f07ff82f9eaae326e1ce8ca8a834d8f0c40febf0}).
The finite elliptic family, rational rank-one curves, partial
inversion, rational approximation, positive Schur complement,
and spanning of derivatives are
\srcid{FiniteAbsLeOneMulMemRange.lean}
(\srcid{a525213fbc51d8331f5cf36529a36e5f4385887a}).
The final polynomial argument for a finite symplectic family is
\srcid{ProdCompEqOfFiniteFamily.lean}
(\srcid{0d643fd13242090c4d63d7543da06e3fab784482}).
The exact change of theta normal form is
\srcid{ThetaLawSymplecticNormalForm.lean}
(\srcid{6950e84acb9e8fa4fc102ec0cb3baef827e74e78}),
and the rescaling and finite-characteristic-family application
are \srcid{ExistsRescale.lean}
(\srcid{193bbb1a07dfd9582c3d8692d5710c88dabb4500}).
The realization of a testing key datum at each period is
\srcid{ExistsKeyInNormalAt.lean}
(\srcid{189fd3b621d2d5c557c11140849fa5f549a6b81c}),
with its explicit \(\mathrm{TorusAlgebraic}\) hypothesis.
